% =========================================================
% Compléments M1 — Surface d'IV (C. Dérivations)
% =========================================================

\subsection*{(1) Breeden--Litzenberger détaillé : du call à la densité}
On part du prix d'un call européen (taux constants pour simplifier) :
\[
C(K,T)=e^{-rT}\E^{\Qmeas}[(S_T-K)^+].
\]
Écrivons cette espérance sous forme intégrale, avec densité $f_{S_T}^{\Qmeas}$ :
\[
C(K,T)=e^{-rT}\int_0^{\infty}(s-K)^+\,f_{S_T}^{\Qmeas}(s)\,\dd s
=e^{-rT}\int_K^{\infty}(s-K)\,f_{S_T}^{\Qmeas}(s)\,\dd s.
\]
\paragraph{Première dérivée en $K$.}
On dérive (formellement, sous conditions de régularité permettant l'échange dérivée/intégrale) :
\[
\frac{\partial C}{\partial K}(K,T)
=e^{-rT}\frac{\partial}{\partial K}\int_K^{\infty}(s-K)\,f(s)\,\dd s,
\quad f=f_{S_T}^{\Qmeas}.
\]
Utilisons la règle de Leibniz :
\[
\frac{\partial}{\partial K}\int_K^{\infty}g(K,s)\,\dd s
=-g(K,K)+\int_K^{\infty}\partial_K g(K,s)\,\dd s,
\quad g(K,s)=(s-K)f(s).
\]
Ici, $g(K,K)=0$ et $\partial_K g(K,s)=-f(s)$. Donc
\[
\frac{\partial C}{\partial K}(K,T)
=e^{-rT}\int_K^{\infty}(-f(s))\,\dd s
=-e^{-rT}\Qmeas(S_T\ge K).
\]
\paragraph{Seconde dérivée.}
On dérive à nouveau :
\[
\frac{\partial^2 C}{\partial K^2}(K,T)
=-e^{-rT}\frac{\partial}{\partial K}\int_K^{\infty} f(s)\,\dd s
=e^{-rT}f(K).
\]
On obtient l'identité :
\[
\boxed{\frac{\partial^2 C}{\partial K^2}(K,T)=e^{-rT}f_{S_T}^{\Qmeas}(K)\ge 0.}
\]
\paragraph{Message.}
La convexité en $K$ équivaut à la positivité de la densité risque-neutre : c'est la base des tests d'arbitrage sur une surface.

\subsection*{(2) Dupire (forme simplifiée) : de $C(K,T)$ à une volatilité locale}
L'idée de Dupire est : \emph{si l'on connaît $C(K,T)$ pour tous $K,T$, on peut fabriquer un modèle de diffusion locale qui reproduit exactement ces prix vanilles}.

Pour garder l'algèbre lisible, on se place dans un cadre simplifié :
\[
r=0,\qquad q=0.
\]
Alors le prix d'un call est simplement
\[
C(K,T)=\E^{\Qmeas}[(S_T-K)^+].
\]
Sous un modèle de volatilité locale
\[
\dd S_t=\sigma_{\mathrm{loc}}(S_t,t)\,S_t\,\dd W_t^{\Qmeas},
\]
la densité $f(s,T)$ de $S_T$ satisfait une équation de Fokker--Planck (équation «forward»).
Dans ce cadre, Dupire montre que $C(K,T)$ satisfait la PDE en variables $(K,T)$ :
\[
\partial_T C(K,T)=\frac12\sigma_{\mathrm{loc}}^2(K,T)\,K^2\,\partial_{KK}C(K,T).
\]
\paragraph{Interprétation.}
Comme $\partial_{KK}C=e^{-rT}f$ (ici $r=0$), la dérivée seconde joue le rôle de densité, et $\partial_T C$ encode comment cette densité «diffuse» dans le temps.

En isolant $\sigma_{\mathrm{loc}}$ :
\[
\boxed{
\sigma_{\mathrm{loc}}^2(K,T)=\frac{\partial_T C(K,T)}{\tfrac12 K^2\,\partial_{KK}C(K,T)}.
}
\]
\paragraph{Remarque.}
Quand $r$ et $q$ sont non nuls, on obtient une formule plus longue (termes en $\partial_K C$ et $C$). Le message reste : la volatilité locale est un \emph{objet dérivé} de $C(K,T)$, donc très sensible au bruit et à l'interpolation.

\subsection*{(3) Contrôles discrets : convexité et densité via différences finies}
En pratique, on observe $C(K_i,T)$ sur une grille discrète de strikes.
Un test simple de convexité (butterfly arbitrage) utilise des différences secondes :
\[
\Delta^2 C_i \approx \frac{C(K_{i+1},T)-2C(K_i,T)+C(K_{i-1},T)}{(\Delta K)^2}.
\]
La condition «pas d'arbitrage butterfly» impose $\Delta^2 C_i\ge 0$ (à tolérance près).
Si $\Delta^2 C_i<0$ à de nombreux endroits, l'interpolation ou les données sont incohérentes économiquement.

\subsection*{(4) Mini-exercice corrigé}
\paragraph{Exercice --- densité implicite à partir de 3 strikes.}
On observe $C(K_{i-1},T)$, $C(K_i,T)$, $C(K_{i+1},T)$ avec un pas uniforme $\Delta K$.
Donner une approximation de la densité risque-neutre en $K_i$ et expliquer le signe attendu.

\emph{Solution.}
Par Breeden--Litzenberger, $f(K_i)\approx e^{rT}\partial_{KK}C(K_i,T)$.
On approxime $\partial_{KK}C$ par la différence seconde :
\[
f(K_i)\approx e^{rT}\frac{C(K_{i+1},T)-2C(K_i,T)+C(K_{i-1},T)}{(\Delta K)^2}.
\]
Comme une densité est positive, on attend un résultat $\ge 0$ (à bruit près). Un résultat fortement négatif indique un arbitrage butterfly (ou des erreurs de données/interpolation).
